%---------------------------------------------------------------------------%
%->> Backmatter
%---------------------------------------------------------------------------%
\chapter[致谢]{致\quad 谢}\chaptermark{致\quad 谢}

在本文完成之际，我谨向所有给予我支持与帮助的师长、同学及家人表示由衷的感谢。

首先，特别感谢我的导师。在硕士研究生期间，导师严谨治学的科研态度和深厚的学术造诣让我受益匪浅。从论文选题、方案设计到最终的撰写定稿，每一个环节都凝聚了导师的辛勤心血。他对我研究中遇到的每一个难点都能给予精辟的指点，并在生活上给予了无微不至的关怀。

感谢实验室的各位同窗。在科研工作的无数个日夜里，我们共同探讨算法实现中的细节，互相鼓励，营造了良好的学术氛围。特别感谢在 FPGA 开发板调试过程中给予我技术支持的同学们，你们的帮助是我顺利完成实验的关键。

感谢中国科学院计算机网络信息中心提供的优越科研条件。在这里，我不仅掌握了专业知识，更学会了如何独立面对复杂的工程难题。

最后，深深地感谢我的家人。正是你们多年如一日的默默支持与理解，才让我能够静下心来完成学业。你们的关爱是我前进道路上最坚实的后盾。

\rightline{2026年6月}
\chapter{作者简历及攻读学位期间发表的学术论文与其他相关学术成果}

\section*{作者简历：}
2019年9月——2023年6月，在福州大学物理与信息工程学院获得学士学位。

2023年9月——2026年6月，在中国科学院上海微系统与信息技术研究所攻读硕士学位。

% 工作经历：

\section*{已发表（或正式接受）的学术论文：}

{
\setlist[enumerate]{}% restore default behavior
\begin{enumerate}[nosep]
    \item BU Yiheng, LUO Yihua, WANG Lei, et al. Hardware/Software Co-design and Implementation of FPGA-based SLAM Backend Accelerator[J]. Robot, 2025.
    % \item 已发表的工作2
\end{enumerate}
}

\section*{申请或已获得的专利：}

王磊,卜怡恒,罗艺华,朱冬晨,李嘉茂.基于稀疏矩阵分解的 SLAM 后端优化方法及硬件加速器架构,中国,发明专利,申请号:202610199997.0

% \section*{参加的研究项目及获奖情况：}


\cleardoublepage[plain]% 让文档总是结束于偶数页，可根据需要设定页眉页脚样式，如 [noheaderstyle]
%---------------------------------------------------------------------------%
